\section{The One Idea: Reality as a Crystal of Experience}

Here is the whole theory in one breath, so you know where the staircase leads. Everything after this is detail.

\subsection{The reversal}

The usual picture of the world starts with dead matter. Tiny lifeless particles bump around, and somehow, after billions of years and enough complexity, some of them wake up and start to feel. Nobody can say how that last step works. How does feeling come from stuff that does not feel? This is the famous hard problem, and on the usual picture it is a brick wall.

Vibe Theory turns the picture around. It starts with experience as the ground floor. The smallest thing that exists is not a dead particle but a tiny unit of feeling, called a \textbf{vibe}. Matter is not the foundation. Matter is what experience looks like once it has settled into stable, shared, repeating patterns. A rock is not dead stuff that might someday feel. A rock is a region of the field of experience that has frozen into a very steady, very boring pattern.

On this view your mind is not produced by your brain. Your brain is a particularly intricate, tightly knit knot in the one field of experience. The feeling was never the mystery to be explained away. It was there from the start, and the thing to explain is how the solid, lawful, shared world we all see gets built out of it. That is the debt this theory pays.

\subsection{So is a rock conscious?}

The reversal raises an obvious worry, and it is worth settling immediately so it does not nag at you. If a rock is frozen feeling, is the rock conscious? Is your coffee mug suffering? Is an electron having a tiny inner life?

The answer is that feeling comes in degrees, and almost everything is at the very bottom of the scale. A single vibe is the faintest imaginable flicker, not a mind, more like one dim pixel than a picture. A rock is a huge number of such flickers locked into a dull, barely-coordinated pattern, each doing almost nothing, none of them woven together into a whole. So there is, at most, next to nothing it is like to be a rock, the way there is next to no music in a single held note. Rich experience, the kind you have, is not present wherever there is matter. It requires vast numbers of vibes bound tightly together into a coordinated whole, and most matter is nowhere close. Later we will meet the exact dial for this, called integration, that says how much a clump of vibes adds up to a genuine experiencer. For a rock it is essentially zero. For a brain it is enormous.

So this is not the cartoon view that your furniture has feelings and your toaster might be in pain. It is the more careful claim that feeling is the basic ingredient, present in the faintest possible form at the bottom, and that it only becomes anything like a mind when it is organized, integrated, and alive. The bottom of reality feels in the way a single grain of sand is a beach: technically of the same substance, but not yet the thing we mean by the word.

\subsection{The crystal}

So picture the field of experience as a vast crystal. Not a crystal sitting in space, the crystal is what space is made of. Each tiny cell of the crystal is one vibe, one quantum of feeling. Each vibe carries a simple felt charge, called its \textbf{tone}, which is one of three values: pain, peace, or pleasure. When two cells touch, that touching is one vibe experiencing another, called a \textbf{note}.

That is the entire inventory. Colored cells, and the touching between them. Color is feeling. Touch is relation. There is nothing else.

\begin{figure}[h]
\centering
\includegraphics[width=0.47\textwidth]{vibe-mesh-73.pdf}
\caption{The crystal of experience, drawn in two dimensions: a tiling of the curved hyperbolic plane. Each cell is a vibe, its color its tone (red pain, green peace, blue pleasure), and each shared edge is a note, one vibe experiencing another. This flat picture is the simplest face of the real three-dimensional crystal.}
\end{figure}

The crystal is not a static thing. It grows, adding new cells at its edge, forever. And one simple rule runs over it: each cell keeps resetting its tone based on the tones of the cells it touches. Run that rule on that growing crystal, and, the theory claims, you get everything. Space is the large-scale shape of the crystal. Time is its growth. Matter, light, gravity, atoms, and minds are all patterns in the colored, growing web.

\subsection{Why a crystal, and why curved}

Why a crystal and not just a random pile? As we will earn slowly, the crystal has to live in a curved, saddle-shaped space called hyperbolic space, for two reasons. It needs endless room to grow, and curved space has that. And it must not have any built-in preferred direction, because our universe does not, and a curved crystal, surprisingly, has none.

The simplest version is a flat-looking tiling of seven-sided cells. The real one is a three-dimensional crystal of twelve-sided cells. Both are coming up. For now, hold the image: a glowing, growing, colored honeycomb, where each tile is a feeling and each shared edge is one feeling touching another.

\textbf{Takeaway:} experience comes first, matter comes later as its stable patterns, and the whole of reality is one growing, curved crystal of three-valued feeling, run by one simple rule.
